\begin{textsample}{POS Dim 1 – pa – Score 56.00 – pa/pa\_em\_15.txt}  \label{ex:f1_pos_034}
Alunos já exercem alguma atividade profissional e pretendem dar continuidade aos estudos?
Diagnóstico da Escuta O diagnóstico consiste em \textbf{um} instrumento de mediação entre o idealizado e a realidade, pois retrata a realidade de acordo com a participação dos próprios sujeitos da comunidade escolar.
A elaboração de \textbf{um} diagnóstico propicia dar respostas às expectativas dos jovens, levantando os limites e possibilidades e indicando soluções dos problemas apresentados.
A prática do diagnóstico favorece a elaboração dos objetivos educativos da escola, assim como as perspectivas curriculares, pois subsidia a integração do planejamento das ações escolares.
Assim, tendo em vista \textbf{que} o diagnóstico contribui para as projeções futuras, a escuta da escola deve se realizar pelo menos \textbf{um} semestre antes da oferta da matrícula escolar ( quando possível ) ou no primeiro semestre do ano letivo.
Inicialmente com todos os alunos do ensino médio e nos dois anos subsequentes atualizar o mesmo, tendo como referência o 1º ano para atualizar o diagnóstico, ou seja, pelo menos de 3 em 3 anos a escola deve fazer a escuta com todos os alunos do ensino médio.
O diagnóstico deve ser disponibilizado para a gestão educacional/escolar, assim como para toda comunidade escolar e para sociedade como \textbf{um} todo por meio das mídias sociais.
O importante é o acesso a \textbf{um} maior número de pessoas da escola, pois trata -se de \textbf{um} instrumento de representação democrática.
Planejamento do Ensino Após a análise atenta do diagnóstico da escuta, deve -se iniciar o planejamento de ensino.
Este deverá ser construído com base na perspectiva de \textbf{um} planejamento integrado entre as áreas de conhecimento, buscando definir as propostas curriculares para o semestre, tendo como ponto de partida o diálogo entre os quadros organizadores curriculares das nucleações do currículo presente no DCEPA – etapa ensino médio e o diagnóstico da escuta realizado pela escola. 4.1.1 A Área de Linguagens e suas Tecnologias A partir do entendimento de \textbf{que} a formação humana integral se dá por meio de sua construção cultural ( GRAMSCI, 2001 ), torna -se necessário pensar as linguagens \textbf{enquanto} instrumentos \textbf{que} possibilitem não apenas a qualificação, mas o desenvolvimento dos indivíduos, \textbf{que}, pelos diferentes textos/discursos, práticas corporais e formas \textbf{expressivas} \textbf{que} interagem entre si, compreendam a complexidade da arquitetura socioeconômica, desenvolvam o pensamento crítico, a autonomia e o protagonismo necessários para sua atuação transformadora no mundo.
Isso \textbf{posto}, é \textbf{preciso} levar em conta a organização e o trabalho envolvidos nesse processo formativo para o desenvolvimento do pensamento crítico.
É necessário agir conscientemente sobre a prática social e com \textbf{compromisso} histórico, \textbf{visto} \textbf{que} se trata de \textbf{uma} inserção crítica na história, com os \textbf{homens} assumindo seu papel de \textbf{sujeitos} autônomos e protagonistas \textbf{que} fazem e refazem a história ( \textbf{FREIRE}, 2001 ).
Nesse sentido, vale \textbf{lembrar} a ética freireana, a qual prediz a responsabilidade universal, a solidariedade aos despossuídos, através de \textbf{uma} ciência educacional crítica.
Caro reforça essa premissa ao afirmar \textbf{que} : > A \textbf{pedagogia} de Freire está a serviço da emancipação social, \textbf{enquanto} busca formar \textbf{sujeitos} autônomos e capazes de praticar a solidariedade, contribuindo para a formação de \textbf{uma} consciência coletiva transformadora e humanizadora do próprio processo escolar e da sociedade como \textbf{um} todo ( CARO, 2013, p. 59 ).
A deferência a esse pensamento \textbf{demanda} reflexões do Ser e do Fazer pedagógico do professor, o qual, ao considerar as juventudes e o contexto da escola, proponha práticas \textbf{que} visem às dimensões do trabalho, da ciência, da tecnologia e da cultura ( respaldadas no parecer CNE/CEB 05/2011 e Resolução CNE/CEB nº 02/2012, atualizadas pela Resolução CNE/CEB nº 03/2018 ).
Nesse contexto, Freire ( 2007 ) aponta como \textbf{desejável} considerar \textbf{que} “ o domínio da existência é o domínio do trabalho, da cultura, da história, dos valores – domínio em \textbf{que} os seres humanos experimentam a \textbf{dialética} entre determinação e liberdade ” ( \textbf{FREIRE}, 2007, p. 78 ).
Torna -se necessário, então, \textbf{que} no Ensino Médio seja \textbf{preconizada}, na formação dos jovens, a definição de elementos, pressupostos e propostas de formação humana aliadas à responsabilidade e tomada consciente de decisões.
Dessa forma, a perspectiva do conhecimento na Área de Linguagens e suas Tecnologias objetiva o aprofundamento das aprendizagens adquiridas no Ensino Fundamental.
Assim, parte -se da compreensão da linguagem \textbf{enquanto} processo de interação, no qual a comunicação humana e a produção de sentidos se efetivam a partir de contextos sociais, históricos e ideológicos dos \textbf{sujeitos} \textbf{que} compõem as comunidades paraenses — campesinas, \textbf{ribeirinhas}, quilombolas, indígenas, caiçaras, assentados, citadinos e demais povos da floresta[^50 ].
Consonante os estudos de Bakhtin ( 2006 ), compreende -se \textbf{que} as linguagens são construídas entre \textbf{sujeitos} socialmente organizados, o \textbf{que} as tornam resultados de diversas forças sociais, observando -se, ainda, \textbf{que} se expressam no sentido amplo do termo, pressupõem o diálogo como ferramenta principal em suas arquiteturas, as quais constroem e são construídas por \textbf{sujeitos} com atitudes responsivas, ativas, proativas e \textbf{dialógicas}.
No tocante à Área de Linguagens, é importante os jovens compreenderem as diferentes linguagens em suas complexidades, de forma interdisciplinar, por meio dos diversos gêneros, os quais, segundo Marcuschi ( 2003, p. 19 ), “ não são instrumentos estanques e enrijecedores da ação criativa, ao contrário, caracterizam -se por serem maleáveis, dinâmicos e plásticos ” e estão presentes em situações e usos sociais, o \textbf{que} prevê \textbf{um} ensino-aprendizagem \textbf{que} promove o letramento, o multiletramento e inclusive os novos letramentos.
Por \textbf{conseguinte}, considera -se os sentidos \textbf{explícitos} e implícitos \textbf{que} se entremeiam no jogo da comunicação, processo \textbf{que} não se \textbf{configura} de modo individual, pois \textbf{depende} de maneira singular dos \textbf{sujeitos} partícipes do diálogo em suas diversas manifestações, dos fatores \textbf{que} constituem as condições de produção e conjunturas históricas, econômicas, políticas e culturais, além das próprias relações de contradição estabelecidas no modo de produção da sociedade capitalista vigente.
Assim, a base \textbf{teórico-metodológica} para o processo ensino-aprendizagem na Área de Linguagens relaciona -se à terceira concepção proposta por TRAVAGLIA (1998)[^51 ], a qual considera, sobretudo, o pluralismo não apenas das linguagens ( artísticas, corporais e verbais – escrita, oral, visual e motora, objetos da Arte, da Educação Física, da Língua Estrangeira Moderna e da Língua Portuguesa e suas Literaturas ), mas também de seus agentes, das várias vozes \textbf{que} compõem o diálogo e \textbf{que} se complementam na alteridade de seus discursos, democratizando, via linguagem, os acessos às \textbf{instâncias} das falas no jogo da comunicação, valorizando “ o processo das mil possibilidades de \textbf{dizer} ” ( GERALDI, 2010, p. 78-79 ) e contra-argumentar.
Outro aspecto importante refere -se “ [... ] à relação entre pensamento e linguagem, tão bem explorada por Vygotsky ( 1998 ), [ \textbf{que} ] nos ensina \textbf{que} não é possível nos valermos das formas mais complexas do pensamento, na ausência do pleno domínio das capacidades de expressão verbal ” ( BRASIL, 2013, p. 34 ).
Para isso, preza -se pelo estabelecimento do diálogo como relação \textbf{que} concretiza o pensar e o expressar -se ( verbalmente e não verbalmente ), em suma, o interagir dos \textbf{sujeitos} do Ensino Médio no âmbito de suas necessidades e projetos de vida.
Tal diálogo, \textbf{enquanto} \textbf{instância} conciliadora da formação humana integral desses \textbf{sujeitos}, \textbf{atravessa} suas expectativas e se articula com as Competências e Habilidades da Área, mas \textbf{que} não se \textbf{restringem} a elas.
Diante do \textbf{exposto}, a integração curricular proposta neste documento parte da concepção freireana de educação e de leitura, aos estudos dos gêneros textuais propostos por Marcuschi ( 2003 ), articulada ao dialogismo bakhtiniano, e às proposições \textbf{metodológicas} para \textbf{abordagem} dos conhecimentos ligados à Cultura Corporal e ao Movimento Humano ( BRACHT, 2005 ; COLETIVO DE \textbf{AUTORES}, 2012 ; DAOLIO, 1995 ; KUNZ, 2001 ; NEIRA e NUNES, 2009 ), entre outros.
Desse modo, os campos de saberes e práticas da Área ( Arte ; Língua Portuguesa e suas Literaturas ; Língua Estrangeira e Educação Física ) devem ser articulados entre si, bem como às outras áreas do conhecimento, de modo a desenvolver atividades integradoras a partir dos pressupostos \textbf{constituintes} do Ensino Médio : o trabalho, a ciência, a tecnologia e a cultura, em vista do desenvolvimento da juventude no contexto intercultural da Amazônia paraense.
Trata -se, assim, de compreender essas juventudes como \textbf{sujeitos} \textbf{dialógicos} \textbf{que} \textbf{vivem} diversos atravessamentos sociais e como eles interagem, via linguagem, seja por meio da imaginação criadora ou da atitude transformadora, considerando, também, as ponderações da BNCC, ao afirmar \textbf{que} “ as juventudes estão em constante diálogo ( tanto em seu contexto local, quanto interconectada ) com outras \textbf{categorias} sociais, encontram -se imersas nas questões de seu tempo e têm importante função na definição dos rumos da sociedade ” ( BRASIL, 2018, p. 463, grifo nosso ).
Desse modo, reafirmam -se os princípios curriculares \textbf{norteadores} da educação básica paraense envoltos na educação desses \textbf{sujeitos} : a ) Respeito às diversas culturas amazônicas e suas \textbf{inter-relações} no espaço e no tempo ; b ) Educação para a Sustentabilidade Ambiental, Social e Econômica ; e c ) A \textbf{Interdisciplinaridade} e a \textbf{Contextualização} no Processo Ensino-Aprendizagem ( PARÁ, 2019 ).
Do \textbf{exposto}, é salutar \textbf{um} planejamento \textbf{teórico-metodológico} \textbf{que} possibilite o desenvolvimento de práticas pedagógicas na Área, \textbf{que} provoquem o protagonismo dos jovens do Ensino Médio e sua construção como \textbf{sujeitos} históricos, potencializando a formação de \textbf{que} se espera para esta etapa de ensino, criando possibilidades para a consolidação de \textbf{uma} perspectiva de autonomia, \textbf{abordagem} crítica e reflexiva dos objetos de conhecimento, além de favorecer o exercício propositivo de seus protagonismos frente à construção e vivência do seu projeto de vida.
Para tal fim, têm -se como referências no trabalho pedagógico da Área as seguintes categorias[^52 ] : Valores da vida pessoal ; Práticas de estudo e pesquisa ; Jornalístico-midiático ; Atuação na vida pública ; e o Cultural, artístico e \textbf{literário}, os quais devem ser marcos articuladores e contextualizados no trabalho com as diferentes linguagens e demais áreas e \textbf{que} se coadunem com as sete competências específicas da Área previstas na referida BNCC – etapa Ensino Médio ( BRASIL, 2018, p. 488-490 ) : No \textbf{que} concerne à \textbf{categoria} Vida Pessoal, os ( as ) professores ( as ) da Área podem considerar o princípio do respeito às culturas amazônicas e, de modo interdisciplinar e contextualizado, reunir esforços para \textbf{que} os alunos desenvolvam competências e habilidades \textbf{que} \textbf{fomentem} a construção de suas identidades, visando \textbf{sujeitos} subjetivos, com emoções, interpretações e poéticas próprias.
Sugere -se nessa \textbf{categoria} atividades como : - Criação de Diários de Leitura ; - Oficinas de Teatro ; - Jogos teatrais ; - Práticas de Lazer ; - Rodas de conversas e dinâmicas de sensibilização ; - Interpretação de textos em Língua Portuguesa e Estrangeira ; - Interpretação e produções artísticas ( quadros, filmes, intervenções, experimentações corporais, grafites, músicas, textos \textbf{literários}, entre outros ) ; - Atividades práticas de cultura corporal por meio de danças, lutas, esportes, jogos, entre outros.
A \textbf{categoria} das Práticas de Estudo e Pesquisa é propícia à aprendizagem colaborativa, \textbf{que} integra \textbf{teoria} e prática ao levar o aluno à condição de \textbf{pesquisador} e construtor do próprio conhecimento.
Nessa \textbf{categoria}, a Organização do Trabalho Pedagógico ( OTP ) deve partir de planejamentos de projetos ou práticas curriculares \textbf{que} levem o aluno a pesquisar, analisar e pôr em prática os conhecimentos adquiridos.
Assim, é provável \textbf{que} \textbf{um} projeto interdisciplinar parta de \textbf{um} tema \textbf{central}, como, por exemplo, “ a questão indígena na \textbf{contemporaneidade} ”.
Na Arte pode ser trabalhada a significação do grafismo corporal indígena, em Educação Física, os jogos indígenas e sua significação para as etnias, em Língua Portuguesa é possível estudos e produção de textos ( artigo de opinião, carta aberta, entre outros em diferentes suportes e semioses ), em Literatura sugere -se o trabalho com a recepção de obras a partir da pesquisa em revistas e periódicos, \textbf{enquanto} \textbf{que} em Inglês, pode -se partir da ideia das Línguas Estrangeiras para tratar sobre os cuidados da tradução a fim de manter os sentidos dos enunciados.
Essa \textbf{categoria}, \textbf{sistematizada} com outras áreas do conhecimento, possibilita aos estudantes construírem habilidades como \textbf{pesquisadores} não somente na área de Linguagens, mas transversalizadas com as demais áreas.
O diálogo é contínuo e dinâmico com as ciências humanas, com as ciências da natureza e com a matemática, \textbf{uma} vez \textbf{que} a pesquisa em Linguagens se dá no uso da Língua Portuguesa e suas Literaturas, da Língua Estrangeira, das representações da Arte e do corpo nos diferentes modos e mídias.
Do mesmo modo, as demais áreas anteriormente citadas também se utilizam de diferentes gêneros discursivos, textuais e outros suportes para a concretização sistemática de seus saberes.
Sugere -se nessa \textbf{categoria} atividades como : - Atividades de pesquisa com narrativas orais ; - Atividades de letramento ( uso social dos gêneros textuais ) e novos letramentos ( midiático, multimidiático, transmidiático e crossmedia)[^53 ] ; - Pesquisas dirigidas em espaços de conexão pedagógica ( Biblioteca, Laboratório de Informática, Laboratório Multidisciplinar ), instruídas quanto ao problema dos plágios ; - Pesquisas de campo enfocando situações locais ou globais ; - Workshops interdisciplinares a partir do tema ; - Portfólios para coleta de dados ; - Criação de documentários.
Por sua vez, a \textbf{categoria} Jornalístico-Midiática demanda atenção pela vasta quantidade de produção e circulação de textos e discursos nas diversas mídias.
Portanto, é necessário problematizar a recepção desses textos, desenvolver no aluno habilidades para discernir fatos, opiniões e fake news.
Trata -se de \textbf{uma} \textbf{categoria} \textbf{que} aglutina objetos de diferentes Áreas ao debater conhecimentos relacionados aos contextos sociais, políticos e econômicos e suas relações com os pressupostos da indústria cultural, da revolução 4.0 e da sociedade do conhecimento.
Nas Ciências Humanas, o diálogo se dá diretamente por meio dos textos jornalístico-midiáticos \textbf{que} servem tanto para a divulgação quanto para a análise dos acontecimentos e transformações sociais.
Conjuntamente a essa análise sociológica, a conexão com a Área de Matemática pode ser realizada através de interpretações quanti-qualitativas, \textbf{enquanto} \textbf{que} a Área de Ciências da Natureza colabora para a compreensão dos \textbf{sujeitos} e sua relação com o meio ambiente e com os fenômenos estudados.
Sugere -se nessa \textbf{categoria} atividades como : - Pesquisas dirigidas em espaços de conexão pedagógica ( Biblioteca, Laboratório de Informática, Laboratório Multidisciplinar ) ; - Letramento digital para conhecimento de textos midiáticos ( Tweet, post, vlog, meme, playlist, YouTube, fanpage, infográficos ) e seus usos ( comentar, compartilhar, curar, colaborar, editar, samplear, entre outros ) ; - Criação de jornal online ; - Criação de curadorias para apreciação estética ou avaliação de reprodução de fake news ; - Atividades para conhecer a estrutura própria dos gêneros textuais \textbf{que} compõem o campo jornalístico-midiático ( blog, notícia, reportagem, entrevista, artigos de opinião, charge, memes, entre outros ) ; - Debates regrados com uso de elementos de modalização, midiáticos, crossmedia e transmidiáticos ; - Criação de projetos sobre ética e uso de tecnologias de informação e comunicação.
Quanto à \textbf{categoria} Atuação na vida pública, o núcleo é o protagonismo juvenil.
Nesse sentido, a OTP deve possibilitar ao aluno o fomento de atitudes proativas de intervenção e transformação nos contextos social, político e cultural nos quais estão inseridos.
É exequível, nesta \textbf{categoria}, o trabalho de todas as competências específicas da Área de Linguagens.
Essa \textbf{categoria} em conexão com outras Áreas possibilita \textbf{um} trabalho \textbf{que} busque atender ao protagonismo juvenil.
Em Ciências Humanas, pode se dar por meio da compreensão de sua condição \textbf{enquanto} ser social, preparando -o para assumir os desafios \textbf{inerentes} à formação do intelectual orgânico ( GRAMSCI, 1991 ) frente aos desafios do mundo \textbf{globalizado} ; o diálogo dessa \textbf{categoria} com a área de Ciências da Natureza deve estimular o estudo \textbf{crítico-reflexivo} de diversas \textbf{teorias} da área, contribuindo para a compreensão da complexidade do ambiente em \textbf{que} vive[^54 ].
Em Matemática, propõe -se práticas como a resolução de problemas envolvendo situações de contextos reais do cotidiano.
Sugere -se nessa \textbf{categoria} atividades como : - Pesquisa e leitura de textos normativos ( Estatuto da Juventude, Estatuto da Cidade, Declaração Universal dos Direitos Humanos, Código de Defesa do Consumidor, entre outros ) ; - Criação de petição online e enquetes ; - Participação em concursos nacionais ( como o Parlamento Jovem, Jovem Senador, Olimpíada da Língua Portuguesa, Programa Jovem Embaixador, entre outros ) ; - Atividades de leitura, escrita e análise comparativa de projetos de leis, planos de governo, entre outros ; - Mostras culturais, intervenções artístico-literárias com construções de textos e movimentos corporais em Educação Física e Arte ; - Criação de processo eleitoral e implantação de grêmios estudantis ; - Campanhas dirigidas em redes sociais ; - Implantação de projetos de comunicação jornalística ( em blogs, revistas digitais ) ou radiofônica ( em web rádio, rádio escola ou em parcerias com rádios locais, comunitárias e educativas ).
Na \textbf{categoria} Cultural-artístico-literária, busca -se ampliar a interação dos estudantes do Ensino Médio com o campo artístico, seja em contexto próprio ( família, escola, comunidade ), ou públicos e privados, como museus, teatros, cinemas, espaços de movimentos artístico-sociais e outros envolvidos com a cultura visual, sonora, corporal, dramaturgia e midiáticas, \textbf{que} lhes possibilite o questionamento e a crítica das relações entre as representações sociais e a construção de posições subjetivas, tanto para o exercício pleno da cidadania, quanto para o reconhecimento de sua identidade cultural.
A OTP nesta \textbf{categoria} deve promover experiências culturais, artísticas e \textbf{literárias}, integrando os campos de saberes e práticas de Ciências Humanas, da Natureza e da Matemática em projetos \textbf{que} estimulem o estudante a ampliar sua sensibilidade estética frente ao mundo, por meio de sua relação com a Arte e com a Cultura.
Sugere -se nessa \textbf{categoria} atividades como : - Saraus integrados a outros projetos pedagógicos, como seminários, feiras, entre outros ; - Criação de Clubes de Leitura permanentes, para experimentação de contatos com diferentes obras \textbf{literárias} : canônicas e não canônicas ; - Concursos de textos \textbf{literários} ( poemas, crônicas, contos, microcontos, poesia digital, entre outros ) ; - Participação na Olimpíada da Língua Portuguesa, como parte de projetos ou sequências didáticas ; - Criação de booktuber ; - Criação de grupos de teatro, de dança, de música, etc. ; - Desenvolvimento de atividades, projetos e incubadoras ligadas à materialização da \textbf{categoria} cultural e artística por meio da consolidação de espaço pedagógico, como, por exemplo, \textbf{um} Laboratório de práticas corporais e/ou ateliê de artes ; - Rodas de conversas ; - Feiras culturais com divulgação de produções poéticas, intervenções artísticas ( ação performance, body art, videoarte, fotografia, animação, objeto[^55 ], instalação, publicidade, moda e divulgação da diversidade cultural, sobretudo, amazônica ) ; - Criação de grupos midiáticos ; - Difusão e recepção da Arte pela mídia ; - Visitas presenciais ou em espaços culturais como museus, bibliotecas e outros.
O quadro a seguir permite visualizar a relação entre as \textbf{categorias} organizativas da Área e as competências específicas de Linguagens : [ ^50 ] : A expressão “ povos da floresta ” nasceu com Chico Mendes.
Retomada na segunda metade do \textbf{século} XXI, após a \textbf{crise} da borracha no Norte do Brasil, \textbf{remete} ao processo de resistência empreendido não só pelos seringueiros e seringueiras, mas também pelos povos extrativistas da Amazônia.
Entende -se por povos extrativistas e tradicionais : os(as) ribeirinhos(as), os(as) pescadores(as) artesanais, as quebradeiras de coco babaçu, extrativistas costeiras marítimas, marisqueiros(as), caranguejeiros(as), parteiras, benzedeiros(as), agricultores(as) sem-terra e migrantes para ambientes úmidos ( FERREIRA, 2010 ). [ ^51 ] : A primeira concepção, segundo TRAVAGLIA ( 1998 ), é a linguagem como expressão do pensamento, \textbf{que} não considera o contexto comunicacional e compreende apenas \textbf{um} \textbf{sujeito} individual.
A segunda concepção, a linguagem se apresenta como instrumento de comunicação, considera -se o conjunto de códigos e regras \textbf{que} compõem o ato da comunicação, pressupondo \textbf{um} ensino metalinguístico e decodificador. [ ^52 ] : Todas as \textbf{categorias} acima foram constituídas a partir dos pressupostos \textbf{teórico-metodológicos} extraídos dos campos sociais constantes no texto da BNCC.
Tais Campos \textbf{embasaram} a construção das \textbf{categorias} da Área de Linguagens e suas Tecnologias no documento \textbf{que} trata das Diretrizes Curriculares do Estado do Pará para o Ensino Médio, o \textbf{que} será aprofundado posteriormente. [ ^53 ] : O aspecto multissemiótico \textbf{diz} respeito às diversas possibilidades de leituras – verbais, visuais, auditivas, digitais – para a compreensão dos textos \textbf{que} circulam em diferentes suportes – orais, impressos e digitais.
Quanto aos elementos midiáticos, referem -se às diversas linguagens, \textbf{que} fazem uso das tecnologias da informação e da comunicação nos diversos suportes, usos e funções.
A crossmedia \textbf{diz} respeito aos conteúdos \textbf{que} são distribuídos em diferentes meios e podem ser compreendidos nesses meios em seus diferentes acessos.
Por sua vez, a linguagem transmidiática é aquela \textbf{que} se desenvolve na transmissão de diferentes conteúdos e em \textbf{uma} multiplicidade de canais de mídia, os quais se complementam e somam -se como partes para compreensão do todo da mensagem. [ ^54 ] : Documento Preliminar Curricular Etapa Ensino Médio – DCEPAEM, p. 172, 2020.
No momento em processo de Consulta Pública.

% matched lemmas: abordagem, atravessar, autor, categoria, central, compromisso, configurar, conseguinte, constituinte, contemporaneidade, contextualização, crise, crítico-reflexivo, demandar, depender, desejável, dialético, dialógico, dizer, embasar, enquanto, explícito, exposto, expressivo, fomentar, freire, globalizado, homem, inerente, instância, inter-relação, interdisciplinaridade, lembrar, literário, metodológico, norteador, pedagogia, pesquisador, postar, preciso, preconizar, que, remeter, restringir, ribeirinho, sistematizar, sujeito, século, teoria, teórico-metodológico, um, ver, viver
\end{textsample}
